
\section{Introdução}

\noindent \begin{minipage}[c]{0.6\textwidth}
  \vspace {1cm}
  \begin{description}
    \item [Banana] Exemplo de mini página com figura e seus respectivos rotulos, para que sejam referenciados ao decorrer do texto.
    \item [Maça] Veja que a Figura \ref{fig:place}, está reservando um espaço para adição de figuras, e o mesmo já esta referenciando seu autor e sua nomeclatura com o indice automatico.
  \end{description}

\end{minipage}
\begin{minipage}[c]{0.4\textwidth}
  \captionof{figure}{Placeholder.}
  \includegraphics[width=\textwidth]{figure/placeholder.jpg}
  \label{fig:place}
  {\fontsize{10pt}{\baselineskip}\selectfont
    Fonte: \citeonline{linux:2023}}
\end{minipage}


\section{Métodos}



\subsection{Resultados}
\lipsum[2-2]

\section{Conclusões}



$\underline{{\color{Blue} }}{\color{red} \, \: \: \left ( \iint_{\widetilde{ \theta }}^{\phi } \right )}  $



\par Aqui vai um exemplo de código em \LaTeX2e.


\begin{figure}[H] %Figuras da aula pratica 1.1
  \center
  \caption{Resultado da aula prática 1.2}
  \subfigure[ Algoritmo.\label{fig:pri2}]{\includegraphics[scale=0.4]{figure/placeholder.jpg}}
  \subfigure[Comportamento.\label{fig:seg2}]{\includegraphics[scale=.4]{figure/placeholder.jpg}}
  % Deixar o espço
  
  {\fontsize{10pt}{\baselineskip}\selectfont Fonte: \citeonline{alguem2022nada}}
  \label{fig:ap1_cod_vigual1}
\end{figure}



Para referenciar utilize \cite{ninguem2022curioso}. Também pode ser citado integrada ao texto, de acordo com \citeonline{alguem2022nada}.

Para inserir imagens adicione a figura no diretório \textit{/figure}


\par Estou usando \href {https://cocalc.com/} {CoCal}

E para referenciar a figura \ref{fig:imagem_massa} utilize dessa forma.

\begin{equation}
  a + b = c \label{1}
  $$ $$ %usado para pular linha
  \sqrt{x} + \sqrt{\smash[b]{y}} + \sqrt{z}
\end{equation}


\par Exemplo de incerção de formula \eqref{1},


Aqui é um exemplo de rodapé. \footnote{Um exemplo de rodapé}


\begin{enumerate}[label=\Roman{*}, ref=(\roman{*})]
  \item fsfsdf
  \item kugfhiuh
\end{enumerate}

\begin{asparaenum}
\item Anterior ... \cite{ninguem2022curioso}
\item Próximo ... \label{pl1}
\end{asparaenum}


\section{Conclusões}




  %$X \xLongleftarrow[\text{NATAN}]{\text{OGLIARI}} Y $ %COM TEXTO
	% $\uparrow$ %Seta para Cima
	%$\overleftarrow{NATAN}$
